% Teaching geometry from eq:ket-sharing-toy-losses; no experimental data.
% Source gradients are g1=(1,0), g2=(-2,2), not (-2,3).
% All panels use 11.5 mm per coordinate unit; canvas is 169 x 74 mm.
% Arrows show gradients or their combinations d; displacement is -eta*d.
\begin{tikzpicture}[
  x=1mm,y=1mm,
  font=\figurefont\footnotesize,
  text=BookInk,
  >={Latex[length=1.7mm,width=1.2mm]},
  axis/.style={draw=BookMuted,line width=0.6pt,-{Latex}},
  tick/.style={draw=BookMuted,line width=0.5pt},
  original/.style={draw=FigureInput,line width=0.9pt,-{Latex}},
  projected/.style={draw=FigureModel,line width=1pt,-{Latex}},
  result/.style={draw=FigureLoss,line width=1.1pt,-{Latex}},
  guide/.style={draw=BookMuted,dashed,line width=0.7pt},
  leader/.style={draw=BookMuted,line width=0.5pt},
  heading/.style={anchor=west,font=\figurefont\footnotesize\bfseries},
  label/.style={inner sep=0.6mm},
  note/.style={align=center,inner sep=0pt}
]
  \path[use as bounding box] (0,0) rectangle (169,74);
  \node[heading] at (0,70) {(a) PCGrad：投影};
  \node[heading] at (58,70) {(b) MGDA：凸包最短点};
  \node[heading] at (116,70) {(c) 单位方向与共同下降};

  % Identical origins, ranges and physical scales in all three panels.
  \foreach \offset in {0,58,116} {
    \begin{scope}[shift={(\offset,0)}]
      \begin{scope}[shift={(29,25)},x=11.5mm,y=11.5mm]
        \draw[axis] (-2.25,0) -- (1.45,0);
        \draw[axis] (0,-1.2) -- (0,2.9);
        \foreach \v in {-2,-1,1} {
          \draw[tick] (\v,-0.045) -- (\v,0.045);
          \node[label,anchor=north] at (\v,-0.09) {$\v$};
        }
        \foreach \v in {-1,1,2} {
          \draw[tick] (-0.045,\v) -- (0.045,\v);
          \node[label,anchor=east] at (-0.09,\v) {$\v$};
        }
        \node[label,anchor=north east] at (-0.06,-0.06) {$0$};
      \end{scope}
      \node[label,anchor=west] at (46.5,25) {$d_1$};
      \node[label,anchor=south] at (29,60) {$d_2$};
    \end{scope}
  }

  \begin{scope}[shift={(29,25)},x=11.5mm,y=11.5mm]
    \coordinate (p1) at (0.5,0.5);
    \coordinate (p2) at (0,2);
    \draw[guide] (1,0) -- (p1);
    \draw[guide] (-2,2) -- (p2);
    \draw[original] (0,0) -- (1,0);
    \draw[original] (0,0) -- (-2,2);
    \draw[projected] (0,0) -- (p1);
    \draw[projected] (0,0) -- (p2);
    \draw[result] (0,0) -- (0.5,2.5);
    % The two removed components are orthogonal to the projected vectors.
    \draw[leader] (0.42,0.42) -- (0.50,0.34) -- (0.58,0.42);
    \draw[leader] (-0.14,2) -- (-0.14,1.86) -- (0,1.86);
  \end{scope}
  \node[label,anchor=west] at (42,29) {$\bm g_1$};
  \node[label,anchor=south] at (6,49) {$\bm g_2$};
  \node[label,anchor=west] at (38,35) {$\widetilde{\bm g}_1$};
  \draw[leader] (p1) -- (38,34);
  \node[label,anchor=south] at (21,54) {$\widetilde{\bm g}_2$};
  \draw[leader] (p2) -- (21,54);
  \node[label,anchor=west] at (37,56) {$\bm d_{\mathrm{PC}}$};
  \node[note] at (26.5,4) {$\bm d_{\mathrm{PC}}=(0.5,2.5)^{\top}$};

  \begin{scope}[shift={(87,25)},x=11.5mm,y=11.5mm]
    \coordinate (mgda) at ({4/13},{6/13});
    \draw[draw=BookInk,line width=1pt] (-2,2) -- (1,0);
    \fill[FigureInput] (-2,2) circle[radius=0.9mm];
    \fill[FigureInput] (1,0) circle[radius=0.9mm];
    \draw[projected] (0,0) -- (mgda);
    \fill[FigureModel] (mgda) circle[radius=0.7mm];
    % Perpendicularity to g1-g2=(3,-2) identifies the closest point.
    \draw[leader]
      ({4/13+0.16*3/sqrt(13)},{6/13-0.16*2/sqrt(13)})
      -- ({4/13+0.16/sqrt(13)},{6/13-0.16*5/sqrt(13)})
      -- ({4/13-0.16*2/sqrt(13)},{6/13-0.16*3/sqrt(13)});
  \end{scope}
  \node[label,anchor=south] at (64,49) {$\bm g_2$};
  \node[label,anchor=west] at (99,29) {$\bm g_1$};
  \node[label,anchor=west] at (96,37) {$\bm d^\star$};
  \draw[leader] (mgda) -- (96,36);
  \node[note] at (84.5,4) {$\bm d^\star=(4/13,6/13)^{\top}$};

  \begin{scope}[shift={(145,25)},x=11.5mm,y=11.5mm]
    % g1.d>0 and g2.d>0 give d1>0 and d2>d1: angles 45 to 90 deg.
    \fill[FigureModel!12] (0,0) -- (45:1)
      arc[start angle=45,end angle=90,radius=11.5mm] -- cycle;
    \draw[draw=BookMuted,dotted,line width=0.7pt] (0,0) circle[radius=11.5mm];
    \draw[guide] (45:1) -- (0,0) -- (90:1);
    \coordinate (unitm) at ({2/sqrt(13)},{3/sqrt(13)});
    \coordinate (unitn) at ({sin(22.5)},{cos(22.5)});
    \draw[projected] (0,0) -- (unitm);
    \draw[result,dashed] (0,0) -- (unitn);
  \end{scope}
  \node[note,anchor=north west,align=left] at (118,54)
    {共同下降区域\\$d_1>0$\\$d_2>d_1$};
  \node[label,anchor=south] at (151,44) {$\bm d_{\mathrm N}$};
  \draw[leader] (unitn) -- (151,44);
  \node[label,anchor=west] at (159,38) {$\bm d_{\mathrm M}$};
  \draw[leader] (unitm) -- (159,38);
  \node[label,anchor=east] at (133,12) {单位圆};
  \draw[leader] (134,13) -- (138.1,15.8);
  \node[note] at (142.5,4) {实际位移：$-\eta\bm d$};
\end{tikzpicture}
